\section{Az alakalmazás további fejlesztési lehetőségei}

\subsection{Szerveres adattovábbítás és elemzés}

A jelenlegi alkalmazás a gyakorlat videóit és a  mérési adatokat a készüléken lokálisan tárolja, az
analízist a telefon saját erőforrásain végzi el és így jeleníti meg az adatokat.
Ez nagy mértékben korlátozza a számítás komplexitását és a gyakorlatok során gyűjtött adatok mennyiségét.
Erre lenne megoldás egy szerveres adattovábbítási és elemzési mechanizmus kialakítása, ez nem
csak a telefont tehermentesítené, de lehetővé tenné olyan matematikai adatelemzési módszerek
alkalmazását, amelyeket más környezet, például Python, R vagy MATLAB keretében könnyebben
lehet megvalósítani, mint Androidon.


\subsection{Platform kompatibilitás}

Ahogy a dolgozat elején is említésre került, a jelenlegi alkalmazás Android rendszerre van
megírva és tesztelve, mivel a Multiplatformos megközelítés lehetővé teszi, hogy később ez az
alkalmazás iOS-en is futtatható legyen, megfelelő teszt eszköz biztosításával az alkalmazás
iOS verziója is létrehozható a lokális környezetnek megfelelő videó- és fájlmenedzsment valamint
a Bluetooth-kapcsolat kezelésének megvalósításával.

\subsection{További adatelemzési módszerek}

A jelenlegi rendszerben az elemzés elsősorban egyszerű statisztikai mutatókra és ökényesen definiált
küszöbértékekre épül. Ezen módszerek szakemberrel való hitelesítése fontos lenne az
alkalmazás teszt verzióján való túllépéshez.

Illetve mivel a betegek testsúlya és mozgásmintázata eltérő, ezek a merev határértékek nem mindenkinél
egyformán pontosak. Egy idősoros adatokra felkészített neurális hálózat (pl. LSTM) vagy egy felügyelt
tanulással betanított modell (pl. Random Forest) képes lenne a nyers szenzoradatokból
automatikusan megtanulni az egyéni mintázatokat, ezáltal sokkal robusztusabbá téve a lépésdetektálást
és a Freezing of Gait (FoG) felismerését.

Gépi tanulás segítségével az alkalmazás képessé válhatna előrejelzések készítésére is,
például a jövőbeli esési kockázat megbecsülésére a járás asszimetriájának vagy a
stabilitás hosszú távú romlásának felismerésével. Implementálható lenne a teljes lépésciklus
folytonos fázisosztályozása (sarokleütés, gördülés, elrugaszkodás és lengési fázis
mikromásodperces elkülönítése).